\section*{Introducción}
\phantomsection
\addcontentsline{toc}{section}{Introducción}
\label{v:presentacion}
Este trabajo establece la derivación de las distribuciones de
Bose--Einstein y Fermi--Dirac a partir de las completaciones cuánticas de
estados con trayectoria y memoria, y demuestra su convergencia al régimen
diluido de Maxwell--Boltzmann. La misma estadística bosónica, realizada
sobre los modos transversales del campo electromagnético, determina la
ley espectral de Planck, el flujo de Stefan--Boltzmann y los
desplazamientos de Wien. El problema central es relacionar la estructura
del estado, sus ocupaciones admisibles y sus observables térmicos mediante
operadores y medidas explícitos. Así, las leyes de ocupación y de
radiación se obtienen en etapas conectadas de una construcción, con las
condiciones de equilibrio y de representación que corresponden a cada una.

La base es la Holografía Modular Triádica (HMT), junto con su realización
dimensional en Mecánica Dimensional. La construcción de caminos
\emph{All Possible Paths} (APP) se organiza sobre el soporte
\((\mathbb Z/9\mathbb Z)^2\). Las traslaciones por
\(\pm(1,0)\) y \(\pm(0,1)\) definen su grafo de Cayley aditivo;
las celdas cuadradas proporcionan una estructura toroidal bidimensional.
Sobre ese mismo soporte se evalúan la suma y el producto, conservando
residuos y cocientes. La firma ternaria orientada, TRIT, toma valores
en \(\{-1,0,+1\}\) y distingue orientación y régimen. Su levantamiento
entero conserva conjuntamente el residuo y el acarreo. El núcleo
topológico de fase, \emph{Topological Phase Kernel}
(TPK), compone selección de fase, transporte, actualización y
prolongación del estado. Su registro retiene la trayectoria y la memoria:
el retorno de fase no reinicia el estado.

Las condiciones iniciales y el refinamiento compatible producen las
secciones genealógicas utilizadas en la realización lineal. La
generación regional, la aritmética de historias, el registro de doce fases,
la determinación de alfa y la incidencia excepcional fijan los
antecedentes de los coeficientes y de sus representaciones. Sobre ese
dominio se construyen el espacio de estados, la evolución unitaria y la
ocupación multiparticular. La compatibilidad con las restricciones de
profundidad precede a su transporte entre realizaciones; en la instancia
finita, el calendario proporciona un Hamiltoniano, un propagador y un
funcional de acción explícitos.

El paso decisivo hacia la estadística es el carácter composicional de
paridad. El sector par se completa mediante potencias simétricas, con
ocupaciones enteras no negativas; el sector impar, mediante potencias
exteriores, cuyos operadores de creación son nilpotentes y cuyos
operadores de ocupación son proyectores. De este modo se obtiene la
exclusión algebraica de Pauli. La identificación relativista de esa
paridad con el espín conserva, además, la compatibilidad de la
representación de rotaciones y las condiciones de localidad. Las multiplicidades determinan los recuentos microcanónicos, y la
factorización de las trazas de Gibbs produce las distribuciones de
Bose--Einstein y Fermi--Dirac. En el régimen termodinámico, la variación
de la entropía combinatoria bajo restricciones de energía y población
proporciona una segunda caracterización del equilibrio. Para el sector bosónico se mantiene la
condición sobre el potencial químico que garantiza la convergencia.
Las cotas del régimen diluido cuantifican el límite común de
Maxwell--Boltzmann.

La realización radiativa aplica la ocupación bosónica a una cavidad
periódica tridimensional con dos polarizaciones transversales, dispersión
lineal y potencial químico nulo. El modo constante se fija como dato de
fondo antes de formar el estado de excitaciones. La densidad de modos y
la energía por excitación producen el espectro de Planck. La convergencia
desde las sumas discretas se demuestra controlando uniformemente el
entorno de frecuencia nula y la cola de altas frecuencias. Los momentos
espectrales dan las densidades de número, energía y entropía; la
integración angular da el flujo de Stefan--Boltzmann. El cambio entre
frecuencia y longitud de onda transforma también la medida mediante su
jacobiano, por lo que los dos máximos de Wien satisfacen ecuaciones
distintas. Cada constante de desplazamiento queda caracterizada por el
máximo de la densidad correspondiente.

La relación entre acción, calendario y temperatura determina las
normalizaciones dimensionales de esa realización. La transducción
térmica distingue la energía de decisión, la información adimensional y
la pareja Boltzmann--temperatura. Su composición con la velocidad
constitutiva enlaza las escalas temporal, energética y espacial de la
radiación. Las hipótesis de continuidad, dominio del generador,
contractividad y convergencia permanecen asociadas a las operaciones
que las requieren.

Entre la transducción y las completaciones, la cadena que fija el portador
consta de renovación trítica, doble graduación, memoria energética positiva,
espectro de Williamson, refinamiento por tres, sección marcada, corrientes de
borde y equivalencia térmica centrada. Esta cadena preserva el
origen que desaparece de una densidad normalizada y lo transporta hasta área,
energía libre, gravitación y radiación. Por ello la evaluación completa no se
presenta como coincidencia decimal, sino como salida de operadores, fibras y
mapas tipados.

\Needspace{13\baselineskip}
La interpretación entrópica distingue las configuraciones de ocupación,
la medida sobre historias y el espectro de un operador de densidad.
La frecuencia cronológica de un calendario no es la medida estacionaria
de su envolvente de historias. Esta última se caracteriza
variacionalmente mediante una identidad de divergencia relativa.
En el sector fermiónico cuasilibre conservador del número e invariante
bajo transformaciones de calibre, el correlador de un cuerpo \(0<C<I\)
y su generador modular \(K_1\) se determinan recíprocamente en dimensión
finita. El transporte, la
reducción y el condicionamiento precisan qué información permanece
observable. La recuperación del registro y la identificación orbital de
operadores se articulan con un balance bilateral que conserva
conjuntamente estado y memoria
(sección~\ref{sec:acoplamiento-recuperacion}).

La unidad areal y el entrelazamiento completan esta relación entre
estructura e información. La identidad modular finita expresa las
variaciones de entropía mediante las áreas sectoriales y un déficit de
entropía relativa no negativo que se resta; con \(\Delta_{\rm mod}\ne0\) y
normalización conocida, el área y el Hamiltoniano modular permiten
recuperar las ocupaciones de ambos sectores. El vínculo entre memoria, ocupación y
radiación se establece, por tanto, conservando los mapas que permiten
pasar de un objeto al siguiente, además de sus valores escalares.
El apéndice~\ref{v:integracion} especifica los dominios de las
realizaciones y las condiciones de esas composiciones.

La realización finita sobre \(\Lambda=(\mathbb Z/27\mathbb Z)^2\)
establece cotas exactas de soporte y entropía, con saturaciones por
subgrupos y anuladores. En el dominio dodecafásico, \(\nu_{120}\) y
\(\nu_{270}\) detectan correlaciones de orden que el histograma trítico
elimina. Sobre el dominio primo, \(R_{\rm vac}\) compone la palabra
esturmiana generada, la realización interna de irreducibles y el reloj
\(\log p\). Su identidad modal se entiende mediante sumas simétricas o
medias de Fejér, y sus tablas constituyen testigos numéricos finitos.

\clearpage
\subsection*{Relaciones entre las construcciones y sus resultados}
El cuadro siguiente identifica la contribución de cada bloque al argumento.
El orden de exposición conserva los antecedentes comunes y después distingue
las operaciones específicas: representación, ocupación, condicionamiento y
lectura radiativa. La representación graduada determina las ocupaciones; la realización
espacial y la dispersión determinan su aplicación radiativa.

\begin{table}[H]
\centering
\small
\renewcommand{\arraystretch}{1.2}
\begin{tabularx}{\textwidth}{@{}>{\raggedright\arraybackslash}p{.29\textwidth}>{\raggedright\arraybackslash}X@{}}
\toprule
Construcción y localización & Construcción y función en el argumento \\
\midrule
Núcleo y prolongaciones, §§\ref{sec:nucleo}--\ref{exc:seccion}
& Estados admisibles, rutas orientadas, memoria y refinamiento; salidas
regionales, registro dodecafásico e incidencia. Estos antecedentes conservan
el origen de los coeficientes y de las representaciones posteriores. \\
Realización cuántica y acción,
§§\ref{v:rec:sec:postulados-cuanticos-genealogicos-rev2}--\ref{sec:hmtf2-cristal-temporal}
& Espacio de estados, evolución unitaria y acción del calendario. La
realización fija los operadores que después intervienen en las energías
y en los estados de equilibrio. \\
Memoria y temperatura,
§\ref{v:ant:sec:ii-kb-aut-desarrollo}
& Fibras de lectura, reversibilidad y pareja Boltzmann--temperatura.
Se distinguen información adimensional, escala energética y coordenatización térmica. \\
Capacidad espectral y puente marcado,
§\ref{v:sec:capacidad-espectral-termica}
& Árbol microscópico ponderado, portador graduado, corrientes de borde,
escalas de capacidad y Williamson, y transductor termométrico variacional. \\
Paridad, completaciones de Fock y distribuciones de equilibrio,
§\ref{v:sec:completaciones-algebraicas} y
§§\ref{v:rec:fock_gibbs:section:01}--\ref{v:rec:ocupacion:section:03};
hitos \ref{v:rec:completaciones:statement:02},
\ref{v:rec:fock_gibbs:statement:01},
\ref{v:rec:thm:factorizacion-gran-canonica-rev3} y
\ref{v:rec:thm:limite-clasico-fd-be-rev2}
& Paridad, Fock, transporte por refinamiento y estado de Gibbs.
Las multiplicidades dan los recuentos; las trazas producen las dos
distribuciones de ocupación, cuyo límite diluido se cuantifica, §§\ref{v:rec:ocupacion:section:02}--\ref{v:rec:ocupacion:section:03}. \\
Historias, área y\newline entrelazamiento
& La incidencia areal--volumétrica determina la medida de historias;
la reducción de estados y el operador de área sostienen la lectura
modular de la entropía. Véanse el teorema~\ref{thm:incidencia-hojas-anexo}
y la sección~\ref{v:ant:sec:area}. \\
Incertidumbre de Fourier, memoria ordenada y vacancias primas,
§\ref{v:n19:sec:incertidumbre-fourier} y
§\ref{sec:f051-prime-vacancy-observer}
& La realización finita establece las cotas de soporte y entropía; los
lectores dodecafásicos distinguen orden y marginales; el observador primo
de vacancias resuelve la correlación mediante la identidad modal
simétrica/Fejér, con testigo histórico hasta \(10^8\), reproducción
independiente hasta \(10^6\) y una cota crítica condicionada por la
uniformidad modal. \\
Ocupación bosónica y leyes de radiación térmica,
§\ref{v:sec:radiacion-termica}; proposiciones
\ref{v:prop:ley-espectral}, \ref{v:prop:stefan} y \ref{v:prop:wien};
límite del lema~\ref{v:lem:limite-termodinamico}
& Los modos transversales, la velocidad constitutiva y la estadística
bosónica determinan las sumas térmicas. Las estimaciones de origen y cola
justifican el límite; sus momentos proporcionan densidades, flujo y
máximos espectrales. \\
\bottomrule
\end{tabularx}
\caption{Construcciones anteriores y utilización posterior en el argumento
estadístico y radiativo.}
\label{v:tab:mapa-demostrativo}
\end{table}

La presencia de antecedentes compartidos responde a esta dependencia efectiva.
Cada bloque conserva la información que necesita la operación siguiente;
la coincidencia de una dimensión o de un valor escalar se acompaña del mapa
que permite reutilizarlo. Así, la estadística recibe una representación
graduada y un Hamiltoniano, mientras la radiación recibe además una
realización espacial y su dispersión.
